\documentclass[11pt,a4paper]{article}

% =========================================================
% THE LAST WORD WAR
% Official Game Archive and Implementation Rulebook
% Steward: Victor A. Adeyemi
% Organization: Da-aim Services
% Version: 1.1
% Year: 2026
% =========================================================

\usepackage[T1]{fontenc}
\usepackage[utf8]{inputenc}
\usepackage{geometry}
\usepackage{enumitem}
\usepackage{booktabs}
\usepackage{array}
\usepackage{longtable}
\usepackage{xcolor}
\usepackage{hyperref}

\geometry{margin=1in}

\hypersetup{
    colorlinks=true,
    linkcolor=blue,
    urlcolor=blue,
    citecolor=blue
}

\setlist[itemize]{leftmargin=1.5em}
\setlist[enumerate]{leftmargin=1.7em}

\title{\textbf{The Last Word War\texttrademark}\\
\large Official Game Archive, Rulebook, and Implementation Guide}

\author{
Victor A. Adeyemi\\
Da-aim Services
}

\date{2026}

\begin{document}

\maketitle

\begin{center}
\textit{
The Last Word War\texttrademark\ is stewarded by
Victor A. Adeyemi under the
Da-aim Services research and development ecosystem.
}
\end{center}

\vspace{1em}

\begin{center}
\textit{A turn-based language, reasoning, creativity, and knowledge-chain game.}
\end{center}

\begin{center}
\textbf{Version 1.1}

Official Rulebook and Historical Archive

The Last Word War\texttrademark

Da-aim Services

2026
\end{center}

\vspace{1em}

\tableofcontents

\newpage

% =========================================================
\section{Game Identity}

\subsection{Name}

The official name of the game is:

\begin{center}
\textbf{The Last Word War\texttrademark}
\end{center}

\subsection{Project Philosophy}

The game is founded upon three guiding statements:

\begin{quote}
War is not the goal of this GAME at all.
\end{quote}

\begin{quote}
All it seeks is to promote HUMAN interactions;
\end{quote}

\begin{quote}
Interactions built firmly upon the previous LAST WORD.
\end{quote}

These statements define the gameplay philosophy, development philosophy, and research philosophy of the project.

\subsection{Game Category}

\textbf{The Last Word War\texttrademark} is a turn-based language, reasoning, and conversational game designed for:

\begin{itemize}
    \item digital multiplayer platforms;
    \item classrooms;
    \item study groups;
    \item social gatherings;
    \item language-learning sessions;
    \item debate and communication training;
    \item competitive tournaments;
    \item digital educational platforms.
\end{itemize}

\subsection{Core Concept}

Players construct sentences under a shared time constraint. Each new sentence begins from the last word of the previous player's sentence.

The software coordinates play.

The players determine the quality, interpretation, and culture of gameplay.

% =========================================================
\section{Gameplay Philosophy}

The Last Word War\texttrademark\ is founded upon a simple principle:

\begin{quote}
The software provides the framework for interaction.

The players provide the meaning.
\end{quote}

Rather than attempting to algorithmically govern every aspect of play, the game intentionally preserves human judgment, interpretation, negotiation, and responsibility.

The software facilitates interaction.

Humans create meaning.

\subsection{Computational Facilitation}

The Last Word War\texttrademark\ adopts the principle of Computational Facilitation.

Computing exists to facilitate interaction rather than govern it. The software provides only the computational infrastructure necessary for coordinated play while leaving interpretation, negotiation, and social responsibility to the participants.

\subsection{Time as the Fundamental Constraint}

Time is the central gameplay constraint.

The game provides a shared interaction space and limits each player's opportunity to act. From these two conditions arise pressure, hesitation, creativity, failure, recovery, humour, competition, cooperation, and play.

Timeout is therefore not merely a technical feature. It is one of the philosophical foundations of the game.

% =========================================================
\section{Game Objective}

The objective is to participate in a structured chain of interaction and, where scoring is enabled, to achieve the highest score or win according to the selected endgame format.

Players should:

\begin{itemize}
    \item respond before the timer expires;
    \item begin from the previous player's last word;
    \item contribute to the ongoing interaction;
    \item follow any rules agreed upon before the match begins;
    \item respect the culture of play established by the participants.
\end{itemize}

% =========================================================
\section{Required Participants}

\begin{itemize}
    \item Minimum number of players: \textbf{2}
    \item Maximum number of players: \textbf{Unlimited}
\end{itemize}

For more than two players, a fixed turn order must be established before the game begins.

% =========================================================
\section{Game Setup}

Before play begins, participants may agree on:

\begin{enumerate}
    \item response timer;
    \item scoring method;
    \item strictness of sentence interpretation;
    \item whether challenges are allowed;
    \item endgame format;
    \item permitted themes or special modes;
    \item local house rules.
\end{enumerate}

\subsection{Example Setup}

\begin{itemize}
    \item Timer: 30 seconds per turn
    \item Mode: Standard Chain Mode
    \item Scoring: Submitted turn = +1
    \item Timeout: $-1$
    \item Endgame: First to 10 triggers Final Round
\end{itemize}

% =========================================================
\section{Core Rule}

Each participant must create a sentence beginning from the last word of the previous participant's sentence.

\subsection{Example}

\begin{quote}
\textbf{Player 1:} Knowledge creates opportunity.
\end{quote}

The last word is:

\begin{quote}
\textbf{opportunity}
\end{quote}

The next player may respond:

\begin{quote}
\textbf{Player 2:} Opportunity often comes after preparation.
\end{quote}

The new last word is:

\begin{quote}
\textbf{preparation}
\end{quote}

% =========================================================
\section{Determining the Last Word}

The software may extract the last word from the previous sentence after removing:

\begin{itemize}
    \item punctuation marks;
    \item emojis;
    \item decorative symbols;
    \item trailing spaces;
    \item non-word characters.
\end{itemize}

\subsection{Example}

\begin{quote}
Learning never stops!
\end{quote}

After removing punctuation, the last word is:

\begin{quote}
\textbf{stops}
\end{quote}

\subsection{Hyphenated Words}

Hyphenated words may be treated as one word unless players agree otherwise before the game begins.

Example:

\begin{quote}
Machine-learning improves prediction.
\end{quote}

The last word is:

\begin{quote}
\textbf{prediction}
\end{quote}

% =========================================================
\section{Player Responsibility}

The Last Word War\texttrademark\ intentionally preserves human judgment.

The software does not determine whether a sentence is creative, meaningful, grammatically correct, humorous, strategic, or socially acceptable.

These decisions belong to the players.

Participants may agree before each match on how strictly they wish to evaluate gameplay.

\subsection{Interpretive Freedom}

Players may discover unconventional strategies, exploit rule boundaries, or interpret situations differently.

Such situations are not treated as failures of the game.

They are treated as part of human-governed gameplay.

\subsection{Culture of Play}

The software defines the framework of play.

The players define the culture of play.

% =========================================================
\section{Software Responsibility}

The software is responsible for objective computational tasks, including:

\begin{itemize}
    \item creating rooms;
    \item connecting players;
    \item preserving shared game state;
    \item managing turn order;
    \item recording sentence history;
    \item extracting the current word;
    \item recording scores;
    \item enforcing timeout constraints;
    \item supporting rematches;
    \item recovering interrupted sessions;
    \item maintaining technical reliability;

    \item removing inactive players after three consecutive timeouts;
    \item preventing new players from joining once a match has started.
    
\end{itemize}

The software intentionally avoids deciding the full meaning, honesty, quality, or social acceptability of every human action during gameplay.

\subsection{Player Inactivity Removal}

To preserve the continuity of multiplayer gameplay, a player who accumulates three consecutive timeout penalties shall be automatically removed from the current match.

A timeout is considered consecutive only if the player fails to submit a sentence on each of their successive turns.

Once removed:

\begin{itemize}
    \item the player no longer participates in the current match;
    \item turn order is automatically adjusted;
    \item scores earned before removal remain part of the match record;
    \item the player may participate again only by joining a new match or a subsequent rematch.
\end{itemize}

This rule exists solely to maintain the flow of gameplay and should not be interpreted as a punishment.

\subsection{Closed Match Rule}

A multiplayer match becomes closed once gameplay has officially started.

After the first turn begins:

\begin{itemize}
    \item no additional players may join the current match;
    \item spectators may observe if supported by the platform;
    \item new participants must wait until a rematch or a newly created room begins.
\end{itemize}

This rule preserves fairness by ensuring that all competitors begin under identical conditions.

%%%%%%%%%%%%%%%%%%%%%%%%%%%%%%%%%%%%%%%%%%%%%%%%%%%%%%%%%%%%%%

\subsection{Match Inactivity Termination}

A match continues for as long as meaningful multiplayer interaction remains possible.

The software shall automatically terminate a match under either of the following conditions:

\begin{enumerate}
    \item Only one active player remains after automatic removal of inactive players; or
    \item Every remaining player has become inactive through consecutive timeout penalties.
\end{enumerate}

If only one active player remains, that player is immediately declared the winner.

If no active players remain, the match is declared abandoned due to inactivity.

%%%%%%%%%%%%%%%%%%%%%%%%%%%%%%%%%%%%%%%%%%%%%%%%%%%%%%%%%%%%%%

\subsection{Design Rationale}

The purpose of these rules is to preserve the continuity of multiplayer gameplay while avoiding unnecessary interruptions.

Inactive players should not prevent active participants from continuing the game.

Consequently, inactivity is handled at the player level before it is considered at the match level.

This approach preserves fairness, encourages active participation, and reflects the project's philosophy that the software should coordinate play while allowing human interaction to remain central.

% =========================================================
\section{Continuity Types}

Continuity describes how a player's sentence relates to the previous sentence.

The interpretation of continuity belongs to the players.

\subsection{Literal Continuity}

The required word is used directly in the next sentence.

\begin{quote}
\textbf{Word:} Energy\\
\textbf{Sentence:} Energy can be transferred between systems.
\end{quote}

\subsection{Conceptual Continuity}

The response uses the required word as a concept rather than merely repeating it mechanically.

\begin{quote}
\textbf{Word:} Energy\\
\textbf{Sentence:} Energy conservation is central to classical mechanics.
\end{quote}

\subsection{Advanced Continuity}

The response advances the previous idea logically, scientifically, narratively, or argumentatively.

\begin{quote}
\textbf{Previous sentence:} Energy can be transferred between systems.\\
\textbf{Next sentence:} Systems exchange energy through work, heat, or radiation.
\end{quote}

Players may choose to reward stronger continuity when playing with advanced or educational rules.

% =========================================================
\section{Difficulty Levels}

Difficulty levels define how strictly players agree to interpret the rules.

\subsection{Easy Mode}

Requirements:

\begin{itemize}
    \item begin with the previous last word;
    \item respond before timeout;
    \item no strict thematic requirement.
\end{itemize}

\subsection{Medium Mode}

Requirements:

\begin{itemize}
    \item all Easy Mode requirements;
    \item reasonable continuity with the previous sentence;
    \item basic avoidance of repetitive play.
\end{itemize}

\subsection{Hard Mode}

Requirements:

\begin{itemize}
    \item all Medium Mode requirements;
    \item additional thematic restriction;
    \item stricter player-governed evaluation;
    \item stronger penalty rules if agreed before play.
\end{itemize}

\subsection{Expert Mode}

Requirements:

\begin{itemize}
    \item all Hard Mode requirements;
    \item responses remain within a specialized knowledge domain;
    \item claims should be defensible or factually reasonable;
    \item vague filler sentences may be rejected by players;
    \item technical vocabulary may receive additional points.
\end{itemize}

% =========================================================
\section{Thematic Modes}

The game may be played under a theme. In thematic modes, players agree that each sentence should remain relevant to the chosen domain.

\begin{itemize}
    \item Physics Mode
    \item Mathematics Mode
    \item Biology Mode
    \item Chemistry Mode
    \item History Mode
    \item German Language Mode
    \item Literature Mode
    \item Quantum Computing Mode
    \item Medical Science Mode
    \item Engineering Mode
    \item Story Mode
\end{itemize}

\subsection{Example: Physics Mode}

\begin{quote}
\textbf{Player 1:} Motion depends on the net force acting on an object.
\end{quote}

Last word:

\begin{quote}
\textbf{object}
\end{quote}

\begin{quote}
\textbf{Player 2:} Objects with mass resist changes in motion because of inertia.
\end{quote}

% =========================================================
\section{Scoring System}

The current digital implementation uses a simple scoring system.

\begin{center}
\begin{tabular}{lc}
\toprule
\textbf{Action} & \textbf{Score} \\
\midrule
Submitted turn & +1 \\
Timeout & $-1$ \\
\bottomrule
\end{tabular}
\end{center}

Players may agree to additional scoring rules for advanced, educational, or tournament play.

\subsection{Optional Advanced Scoring}

\begin{center}
\begin{tabular}{lc}
\toprule
\textbf{Action} & \textbf{Optional Score} \\
\midrule
Strong continuity & +1 \\
Creative contribution & +1 \\
Thematic accuracy & +1 \\
Advanced vocabulary or technical use & +1 \\
Invalid sentence by player agreement & $-1$ or $-2$ \\
Failed challenge & $-1$ \\
Successful challenge & +1 \\
\bottomrule
\end{tabular}
\end{center}

\subsection{Recommended Scoring Principle}

The base score should remain simple.

Additional scoring should be used only when players want a stricter or more educational match.

% =========================================================
\section{Continuity Bonuses}

Players may optionally reward long coherent chains.

\begin{center}
\begin{tabular}{lc}
\toprule
\textbf{Coherent Chain Length} & \textbf{Optional Bonus} \\
\midrule
5 turns & +2 \\
10 turns & +5 \\
20 turns & +10 \\
\bottomrule
\end{tabular}
\end{center}

A chain is broken when players agree that continuity has failed, a player times out, or a locally agreed rule is violated.

% =========================================================
\section{Repetition Rule}

To prevent circular gameplay, players may agree that a final word cannot be reused within a recent-window history.

\subsection{Recommended Rule}

A final word should not reappear within the previous \textbf{10 turns}, unless players agree otherwise.

\subsection{Purpose}

This rule prevents loops such as:

\begin{quote}
Knowledge creates opportunity.\\
Opportunity creates knowledge.
\end{quote}

% =========================================================
\section{Restricted and Weak Words}

Some words are too broad and may weaken gameplay.

\subsection{Examples of Weak Words}

\begin{itemize}
    \item thing
    \item stuff
    \item something
    \item anything
    \item everything
    \item good
    \item bad
    \item nice
\end{itemize}

\subsection{Recommended Handling}

Players may choose one of the following options:

\begin{enumerate}
    \item Ban weak words completely.
    \item Allow weak words but give no bonus points.
    \item Penalize repeated weak-word usage.
    \item Accept weak words as part of casual play.
\end{enumerate}

% =========================================================
\section{Word Energy System}

The Word Energy System is an optional advanced scoring mechanic. It assigns strategic value to the final word chosen by each player.

\begin{center}
\begin{tabular}{lc}
\toprule
\textbf{Word Type} & \textbf{Energy Value} \\
\midrule
Common everyday word & 1 \\
Specific noun or verb & 2 \\
Technical or academic term & 3 \\
Rare or specialized term & 4 \\
Interdisciplinary or high-level concept & 5 \\
\bottomrule
\end{tabular}
\end{center}

\subsection{Strategic Purpose}

Players may deliberately end their sentences with difficult words to challenge the next participant. This creates a balance between scoring, creativity, and competitive pressure.

% =========================================================
\section{Challenge Mechanism}

Any player may issue a challenge when a rule violation is suspected.

Challenges are resolved by the participating players or by a mutually accepted human moderator.

The software does not determine challenge outcomes.

\subsection{Grounds for Challenge}

A challenge may be issued for:

\begin{itemize}
    \item weak or false continuity;
    \item failure to use the required word;
    \item thematic violation;
    \item repetition violation;
    \item use of banned words;
    \item fabricated or obviously false technical claims;
    \item locally agreed rule violations.
\end{itemize}

\subsection{Challenge Outcomes}

Players may agree on one of the following outcomes:

\begin{itemize}
    \item the challenged sentence remains accepted;
    \item the sentence is rejected;
    \item the player loses a point;
    \item the challenger gains a point;
    \item the chain resets from the last accepted word;
    \item play continues without penalty.
\end{itemize}

% =========================================================
\section{Endgame Formats}

\subsection{Format A: First to Target Score}

A player reaching a target score triggers the endgame condition.

\subsection{Format B: Final Round}

The first player to reach the target score triggers a final round, allowing other players one final opportunity to respond.

\subsection{Format C: Sudden Death}

Tied leaders may enter Sudden Death until a winner emerges or a draw condition is reached.

\subsection{Format D: Highest Score Wins}

The game ends when the agreed duration expires. The participant with the highest score is declared the winner.

\subsection{Format E: Team Battle}

Players are divided into teams. Each team collaborates internally, but only the assigned player may submit the official sentence during a turn.

\subsection{Format F: Classroom Assessment}

The instructor selects a theme and evaluates students based on agreed learning objectives.

% =========================================================
\section{Special Modes}

\subsection{Blitz Mode}

Response timer: \textbf{30 seconds}

\subsection{Lightning Mode}

Response timer: \textbf{10 seconds}

\subsection{Story Mode}

Every sentence must contribute to a single evolving narrative.

\subsection{Reverse Mode}

Players use the first word instead of the last word.

\subsection{No-Repetition Mode}

No final word may be repeated during the entire round.

\subsection{Expert Citation Mode}

For academic or scientific themes, players may be required to make factually defensible statements. In formal settings, a human moderator may request a short explanation.

% =========================================================
\section{Attack and Defense Cards}

Attack and Defense Cards are optional mechanics for competitive versions of the game.

\subsection{Attack Card Examples}

\begin{itemize}
    \item \textbf{Theme Shift:} The next player must answer within a new theme.
    \item \textbf{Technical Term:} The next player must use a technical concept.
    \item \textbf{No Common Words:} The next player may not end with a common word.
    \item \textbf{Double Timer Pressure:} The next player has half the normal time.
    \item \textbf{Definition Demand:} The next player must define the required word in their sentence.
\end{itemize}

\subsection{Defense Card Examples}

\begin{itemize}
    \item \textbf{Extra Time:} Gain 30 additional seconds.
    \item \textbf{Challenge Shield:} Cancel one challenge.
    \item \textbf{Word Swap:} Replace the required word with a synonym.
    \item \textbf{Theme Immunity:} Ignore one theme-shift attack.
\end{itemize}

% =========================================================
\section{Official Scoreboard Template}

\begin{verbatim}
THE LAST WORD WAR

Mode: Standard
Timer: 30 Seconds
Theme: General

Oasis .......... 12
Victor ........... 10
Ahmed ........... 8
Grace ........... 7

Current Word:
OPPORTUNITY

Current Turn:
Ahmed

Recent Sentence History:
Science helps us understand nature.
Nature inspires technological innovation.
Innovation drives economic growth.

=================================
\end{verbatim}

% =========================================================
\section{Implementation Logic}

This section describes the digital implementation philosophy.

\subsection{Core Data Structures}

The software may track:

\begin{itemize}
    \item player names;
    \item player scores;
    \item current turn index;
    \item previous sentence;
    \item current required word;
    \item sentence history;
    \item timer status;
    \item game phase;
    \item room state.
\end{itemize}

\subsection{Turn Algorithm}

\begin{enumerate}
    \item Display the current required word.
    \item Start or synchronize the timer.
    \item Receive player sentence.
    \item If timeout occurs, apply timeout penalty. If the current player reaches three consecutive timeouts, remove the player from the match and update the turn order.
    
    \item Extract the final word from the submitted sentence.
    \item Record the sentence.
    \item Update the score.
    \item Update the current word.
    \item Advance to the next player.
    \item Persist the shared game state.
\end{enumerate}



\subsection{Pseudocode}

\begin{verbatim}
initialize room
initialize players
initialize scores
set starting word

while game_not_over:
    current_player = players[current_turn]
    display current_word
    synchronize timer

    sentence = receive_input(current_player)

    if timeout:
        scores[current_player] -= 1
        advance_turn()
        persist_state()
        continue

    
    record_sentence(sentence)
    current_word = extract_last_word(sentence)
    scores[current_player] += 1

    evaluate_endgame_state()
    advance_turn()
    persist_state()

announce_result()
\end{verbatim}

If the match has started,
reject all requests to join the room as an active player.

The algorithm coordinates play but does not judge the full meaning or quality of player contributions.

%%%%%%%%%%%%%%%%%%%%%%%%%%%%%%%%%%%%%%%%%%%%%%%%%%%%%%%%%%%%%%

\section{Multiplayer Integrity Rules}

The following rules preserve fairness and continuity during online multiplayer gameplay.

These rules are enforced automatically by the software because they concern the coordination of play rather than the interpretation of gameplay.

%%%%%%%%%%%%%%%%%%%%%%%%%%%%%%%%%%%%%%%%%%%%%%%%%%%%%%%%%%%%%%

\subsection{Closed Match Rule}

Once the first turn of a match begins, the match becomes closed.

During an active match:

\begin{itemize}
    \item no additional players may join the current game;
    \item players already participating continue until the match ends or they are removed through the inactivity rule;
    \item new participants must wait until a rematch or a newly created room begins.
\end{itemize}

This rule ensures that every participant begins the match under the same conditions.

%%%%%%%%%%%%%%%%%%%%%%%%%%%%%%%%%%%%%%%%%%%%%%%%%%%%%%%%%%%%%%

\subsection{Player Inactivity Rule}

A player who fails to submit a sentence before the timeout receives a timeout penalty.

If the same player accumulates three consecutive timeout penalties, that player is considered inactive.

%%%%%%%%%%%%%%%%%%%%%%%%%%%%%%%%%%%%%%%%%%%%%%%%%%%%%%%%%%%%%%

\subsection{Automatic Player Removal}

When a player becomes inactive through three consecutive timeout penalties, the software automatically removes the player from the current match.

Upon removal:

\begin{itemize}
    \item the player's existing score remains recorded;
    \item the player's previous sentences remain part of the sentence history;
    \item the remaining players continue the match;
    \item turn order is automatically recalculated.
\end{itemize}

If only one player remains after automatic removals, that player is declared the winner.

%%%%%%%%%%%%%%%%%%%%%%%%%%%%%%%%%%%%%%%%%%%%%%%%%%%%%%%%%%%%%%

\subsection{Consecutive Timeout Counter}

The inactivity counter applies only to consecutive timeout penalties.

A successful turn immediately resets the player's consecutive timeout counter to zero.

Example:

\begin{center}
\begin{tabular}{c c}
\toprule
Turn Outcome & Consecutive Timeout Count\\
\midrule
Timeout & 1\\
Timeout & 2\\
Successful Turn & 0\\
Timeout & 1\\
Timeout & 2\\
Timeout & Removed\\
\bottomrule
\end{tabular}
\end{center}

%%%%%%%%%%%%%%%%%%%%%%%%%%%%%%%%%%%%%%%%%%%%%%%%%%%%%%%%%%%%%%

\subsection{Design Rationale}

These rules exist solely to preserve the continuity of multiplayer gameplay.

They are computational responsibilities rather than gameplay judgments.

The software coordinates participation.

Human players continue to govern interpretation, creativity, and the culture of play.


% =========================================================
\section{Educational Applications}

The game can be adapted for:

\begin{itemize}
    \item vocabulary development;
    \item language learning;
    \item physics education;
    \item history revision;
    \item scientific communication;
    \item debate training;
    \item critical thinking exercises;
    \item classroom icebreakers;
    \item postgraduate seminar warm-ups;
    \item AI-literacy discussions.
\end{itemize}

\subsection{Classroom Example}

In a physics classroom, the instructor may begin with:

\begin{quote}
Forces cause acceleration.
\end{quote}

Students continue the chain while staying within mechanics, energy, momentum, waves, or another selected topic.

% =========================================================
\section{Digital Product Roadmap}

\subsection{Current Product Direction}

The product direction prioritizes reliable Human--Human interaction over automated judgment.

The system should strengthen multiplayer coordination, onboarding, session recovery, timing, and educational adaptation.

\subsection{Minimum Viable Product}

The first playable digital version should include:

\begin{itemize}
    \item player registration;
    \item turn order;
    \item timer;
    \item sentence input;
    \item automatic last-word extraction;
    \item scoreboard;
    \item game history.
\end{itemize}

\subsection{Intermediate Version}

The intermediate version may include:

\begin{itemize}
    \item improved onboarding;
    \item theme selection;
    \item repetition visibility;
    \item tournament brackets;
    \item player statistics;
    \item classroom dashboard;
    \item spectator support.
\end{itemize}

\subsection{Advanced Version}

The advanced version may include:

\begin{itemize}
    \item mobile application;
    \item voice input;
    \item multilingual play;
    \item educational analytics;
    \item adaptive timer settings;
    \item downloadable classroom materials.
\end{itemize}

% =========================================================
\section{Possible Technology Stack}

\subsection{Current Web Version}

\begin{itemize}
    \item Frontend: React and Vite
    \item Backend: Supabase and PostgreSQL
    \item Real-time play: Supabase Realtime
    \item Deployment: Vercel
    \item Version control: Git and GitHub
\end{itemize}

\subsection{Future Mobile Version}

\begin{itemize}
    \item Flutter or React Native for cross-platform development
    \item Supabase or equivalent backend for authentication and realtime state
    \item Mobile-first timer, input, and notification design
\end{itemize}

\subsection{Offline Classroom Version}

\begin{itemize}
    \item Printable cards
    \item Manual scoreboard
    \item Teacher or student moderator
    \item Optional spreadsheet scoring template
\end{itemize}

% =========================================================
\section{Intellectual Property Archive}

\subsection{Project Steward}

The Last Word War\texttrademark\ is stewarded by:

\begin{center}
\textbf{Victor A. Adeyemi}
\end{center}

under the Da-aim Services research and development ecosystem.

\subsection{Archival Statement}

This document records the formal design, gameplay rules, scoring systems, implementation concepts, educational applications, and historical evolution of \textbf{The Last Word War\texttrademark}.

\subsection{Authorship and AI Policy}

Human contributors remain responsible for:

\begin{itemize}
    \item vision;
    \item interpretation;
    \item governance;
    \item implementation decisions;
    \item publication decisions.
\end{itemize}

Artificial intelligence may assist the project's development, documentation, and research. It does not participate in gameplay and is not considered a player, referee, steward, or author of the project.

\subsection{Copyright Notice}

\begin{center}
\textcopyright\ 2026 Victor A. Adeyemi\\
Da-aim Services\\
\textbf{The Last Word War\texttrademark}\\
All rights reserved.
\end{center}

% =========================================================
\section{Version History}

\begin{center}
\begin{tabular}{lll}
\toprule
\textbf{Version} & \textbf{Date} & \textbf{Description} \\
\midrule
v0.1 & 2026 & Original rulebook draft \\
v0.2 & 2026 & Expanded implementation archive \\
v1.0 & 2026 & Repository synchronization and governance update \\
v1.1 & 2026 & Alignment with Gameplay Philosophy and Computational Facilitation \\
\bottomrule
\end{tabular}
\end{center}

\end{document}

